% Adaptado para IFSP-BRA por Profa. Dra. Ana Paula Müller Giancoli - 27/03/2020

% -- Classe Documento
\documentclass[
	% -- opções da classe memoir --
	hidelinks,
	12pt,				% tamanho da fonte
	openright,			% capítulos começam em pág ímpar (insere página vazia caso preciso)
	oneside,			% para impressão em verso e anverso. Oposto a oneside, twoside
	a4paper,			% tamanho do papel. 
	%normalfigtabnum,
	%pnumromarab,
	% -- opções da classe abntex2 --
	%chapter=TITLE,		% títulos de capítulos convertidos em letras maiúsculas
	%section=TITLE,		% títulos de seções convertidos em letras maiúsculas
	%subsection=TITLE,	% títulos de subseções convertidos em letras maiúsculas
	%subsubsection=TITLE,% títulos de subsubseções convertidos em letras maiúsculas
	% -- opções do pacote babel --
	english,			% idioma adicional para hifenização
	french,				% idioma adicional para hifenização
	spanish,			% idioma adicional para hifenização
	brazil,				% o último idioma é o principal do documento
]{abntex2}

\usepackage{array}
\newcolumntype{P}[1]{>{\centering\arraybackslash}p{#1}}

% --
\usepackage[document]{ragged2e}
% ---
% Pacotes básicos 
% ---
% Usa a fonte Helvet que é parecida com Arial
\usepackage{helvet} 
% define como default.
\renewcommand{\familydefault}{\sfdefault} 
% Selecao de codigos de fonte.
\usepackage[T1]{fontenc}	
% Codificacao do documento (conversão automática dos acentos)
\usepackage[utf8]{inputenc}	
% Usado pela Ficha catalográfica
\usepackage{lastpage}	
% Indenta o primeiro parágrafo de cada seção.
\usepackage{indentfirst}		
\usepackage{graphicx}
\usepackage[table]{xcolor}
\usepackage{float}
\usepackage{lipsum}
\usepackage{tabularx}
\captionstyle{\center} %(para as legendas; use \legend para fonte)
\usepackage{caption}

\usepackage{hyperref}

% ---
% Pacotes de citações
% ---
% Paginas com as citações na bibl
\usepackage[brazilian,hyperpageref]{backref}
% Citações padrão ABNT
\usepackage[alf,abnt-emphasize=bf,abnt-etal-text=emph]{abntex2cite}  
% Adequado para atender a norma ABNT 6023:2018 por Profa. Dra. Ana Paula Müller Giancoli - 27/03/2020
\usepackage{url6023} 



% -- 
% Pacote para moldura
% --
% Adicionado para criar uma moldura ao redor da página
\usepackage{fancybox} 
% incluído a moldura da página
\fancypage{\setlength{\fboxsep}{12pt}\fbox}{} 
\usepackage{wrapfig}


%--

% comando para inserir autor e ano
\newcommand{\citeauthorandyear}[1]{\citeauthoronline{#1} (\citeyear{#1})}
\renewcommand{\ABNTEXchapterfontsize}{\bfseries\large}
% https://www.tablesgenerator.com/

\usepackage{float}
\floatstyle{plaintop} % Coloca caption no topo
\newfloat{quadro}{htbp}{lop}
\floatname{quadro}{Quadro}
\newcommand{\listofquadros}{\listof{quadro}{Lista de Quadros}}
\renewcommand{\baselinestretch}{1.5}

% --
% Orientador e coorientador
% --

\newcommand{\orient}[1]{\def \oorient {#1}}
\newcommand{\imprimirorient}{\oorient}

\newcommand{\coorient}[1]{\def \ocoorient {#1}}
\newcommand{\imprimircoorient}{\ocoorient}

\orient{Nome do Orientador com o título}
\coorient{-}
% --
% Autores
% --
\newcommand{\autorA}[1]{\def \oautorA {#1}}
\newcommand{\imprimirautorA}{\oautorA}
\newcommand{\prontuarioA}[1]{\def \oprontuarioA {#1}}
\newcommand{\imprimirprontuarioA}{\oprontuarioA}

\newcommand{\autorB}[1]{\def \oautorB {#1}}
\newcommand{\imprimirautorB}{\oautorB}
\newcommand{\prontuarioB}[1]{\def \oprontuarioB {#1}}
\newcommand{\imprimirprontuarioB}{\oprontuarioB}

\newcommand{\autorC}[1]{\def \oautorC {#1}}
\newcommand{\imprimirautorC}{\oautorC}
\newcommand{\prontuarioC}[1]{\def \oprontuarioC {#1}}
\newcommand{\imprimirprontuarioC}{\oprontuarioC}

\prontuarioA{BP9999999}
\autorA{Nome do autor A}

\prontuarioB{BP9999999}
\autorB{Nome do autor B}

\prontuarioC{BP9999999}
\autorC{Nome do autor C}

% ----

% --
% Inicio do Documento
% --

\begin{document}
\selectlanguage{brazil}
\thispagestyle{empty} % remove a numeração desta página
\textual

\begin{wrapfigure}[2]{l}{0.30\textwidth}
\includegraphics[scale=.35]{imagens/IFSP-BRA.png} 
\end{wrapfigure}

\centering
\footnotesize\textbf{Tecnologia em Análise e Desenvolvimento de Sistemas} \\	\footnotesize\textbf{Metodologia de Pesquisa Científica e tecnológica - MPCI4}\\

\titulo{Proposta de Trabalho de Conclusão de Curso (TCC)}
\begin{center}
\vspace{1cm}
    \normalsize \textbf{\imprimirtitulo}
\end{center}
			
\begin{flushleft}
    \small \textbf{1. COMPOSIÇÃO DO GRUPO DE TRABALHO}
\end{flushleft}
\begin{center}
	\begin{tabular}{|m{0.5cm}|m{4cm}|m{10cm}|}
		\hline 
	    \textbf{n$^o$.} & \multicolumn{2}{c|}{\textbf{Nome do aluno}} \\
		\hline
		\textbf{1} & \textit{\imprimirprontuarioA} & \textit{\imprimirautorA}\\
		\hline
		\textbf{2} & \textit{\imprimirprontuarioB} & \textit{\imprimirautorB}\\
		\hline
		\textbf{3} & \textit{\imprimirprontuarioC} & \textit{\imprimirautorC}\\
		\hline
	\end{tabular}
\end{center}
	
\begin{flushleft}
   \small \textbf{1.1. PROFESSOR(ES) ORIENTADOR(ES) }
\end{flushleft}

\begin{center}
	\begin{tabular}{|m{5cm}|m{10cm}|}
	    \hline 
	    \multicolumn{2}{|c|}{\textbf{Nome}} \\
		\hline
		\textbf{Orientador} & \textit{\imprimirorient} \\
		\hline
		\textbf{Coorientador} & \textit{\imprimircoorient} \\
		\hline
	\end{tabular}
\end{center}

\begin{flushleft}
   \normalsize\textbf{2.TÍTULO E PROPOSTA DO PROJETO}
\end{flushleft}

\begin{flushleft}
   \small \textbf{Título do Projeto:}
\end{flushleft} 
\normalsize
Título do meu projeto\\

\begin{flushleft}
\textbf{Problema de pesquisa (pergunta):}
\end{flushleft} 
\justify


\hspace{0.5cm} Ao

\begin{flushleft}
\textbf{Justificativa (relevância do estudo):}
\end{flushleft} 
\justify

\hspace{0.5cm} Por meio de citação de autores, contextualizar e justificar o porquê que o seu projeto é importante.

\begin{flushleft}
\textbf{Objetivo Geral:}
\end{flushleft} 


\justify

\hspace{0.5cm} O objetivo Geral é ...

\begin{flushleft}
\textbf{Objetivos Específicos:}
\end{flushleft} 
\justify

Para cumprir com o objetivo geral desta pesquisa, espera-se atingir os
seguintes objetivos específicos:

\begin{itemize}
	\item Objetivo 1;
	\item Objetivo 2;
	\item Objetivo n.
\end{itemize}

\begin{justify}
\textbf{Método de Pesquisa (definição da linha metodológica coerente com o problema:}
\end{justify} 
\subsection{Método de Pesquisa}
Definição da linha metodológica coerente com o problema.
Desenvolver ...


\begin{justify}
\textbf{Metodologia do Desenvolvimento do Sistema (descrição do processo de desenvolvimento, objeto de estudo, tecnologias, requisitos, ferramentas e procedimentos de elaboração e de análise do sistema):}
\end{justify} 
\justify

\hspace{0.5cm} Desenvolver ...

\begin{flushleft}
\textbf{Resultados esperados:}
\end{flushleft} 
\justify
\hspace{0.5cm} Por ..

\begin{flushleft}
\textbf{Cronograma (considerar o 5$^o$. e 6$^o$. módulos:}
\end{flushleft} 
\justify
Considerar os períodos 4, 5 e 6.
O Quadro~\ref{qua:cronograma} resume o cronograma de atividades da proposta de trabalho de conclusão de curso (TCC).


\begin{quadro}[H]
    \begin{center}
	    \caption{Cronograma de execução de atividades por semestre.}
		\label{qua:cronograma}
		\begin{tabular}{|m{3.5cm}|m{.2cm}|m{.2cm}|m{.2cm}|m{.2cm}|m{.2cm}|m{.2cm}|m{.2cm}|m{.2cm}|m{.2cm}|m{.2cm}|m{.2cm}|m{.2cm}|m{.2cm}|m{.2cm}|m{.2cm}|m{.2cm}|m{.2cm}|m{.2cm}|}
			\hline
			\centering
			\cellcolor{gray!50}\textbf{Entregáveis} & \multicolumn{18}{c|}{ \cellcolor{gray!50}\textbf{Atividades a serem realizadas / mês}}\\
			\hline
		    \centering
			\textbf{Anos} & \multicolumn{12}{c|}{\textbf{2022}} & \multicolumn{6}{c|}{\textbf{2023}} \\
			\hline
			\centering
            \cellcolor{gray!50}\textbf{Meses}&\cellcolor{gray!50}\textbf{1}&\cellcolor{gray!50}\textbf{2}&\cellcolor{gray!50}\textbf{3}&\cellcolor{gray!50}\textbf{4}&\cellcolor{gray!50}\textbf{5}&\cellcolor{gray!50}\textbf{6}&\cellcolor{gray!50}\textbf{7}&\cellcolor{gray!50}\textbf{8}&\cellcolor{gray!50}\textbf{9}&\cellcolor{gray!50}\textbf{10}&\cellcolor{gray!50}\textbf{11}&\cellcolor{gray!50}\textbf{12}&\cellcolor{gray!50}\textbf{1}&\cellcolor{gray!50}\textbf{2}&\cellcolor{gray!50}\textbf{3}&\cellcolor{gray!50}\textbf{4} &\cellcolor{gray!50}\textbf{5} &\cellcolor{gray!50}\textbf{6}\\
			\hline
			Formação do Grupo &\cellcolor{gray!50} & & & & & & & & & & & & & & & & & \\
			\hline
			Definição do tema &\cellcolor{gray!50} & & & & & & & & & & & & & & & & & \\
			\hline
			Introdução e problema &\cellcolor{gray!50} & & & & & & & & & & & & & & & & & \\
			\hline
			Justificativa e Objetivos & &\cellcolor{gray!50} & & & & & & & & & & & & & & & & \\
			\hline
			Pesquisa de metodologias & & &\cellcolor{gray!50} &\cellcolor{gray!50} & & & & & & & & & & & & & & \\
			\hline
			Resultados esperados & & & & &\cellcolor{gray!50} & & & & & & & & & & & & &\\
			\hline
			Cronograma e Apresentação & & & & &\cellcolor{gray!50} & & & & & & & & & & & & & \\
			\hline
			Revisão Bibliográfica & & & & & &  \cellcolor{gray!50} &\cellcolor{gray!50} &\cellcolor{gray!50} & \cellcolor{gray!50} & & & & & & & & & \\
			\hline
			Levantamento de Requisitos  & & & & & & &\cellcolor{gray!50} &\cellcolor{gray!50} &\cellcolor{gray!50}& & & & & & & & & \\
			\hline
			Elaborar Diagrama de Casos de uso e Classes  & & & & & & & & &\cellcolor{gray!50} &\cellcolor{gray!50} & & & & & & & & \\
			\hline
			Descrições dos Casos de Uso  & & & & & & & & &\cellcolor{gray!50} &\cellcolor{gray!50} & & & & & & & & \\
			\hline
			Elaboração dos Protótipos em ferramenta específica & & & & & & & & & &\cellcolor{gray!50} &\cellcolor{gray!50} &  \cellcolor{gray!50} & & & & & & \\
			\hline
			Entregas e Revisão & & & & & & & & & & & & & &\cellcolor{gray!50} &\cellcolor{gray!50} &\cellcolor{gray!50} & & \\ 
			\hline
			Desenvolvimento  & & & & & & &\cellcolor{gray!50} &\cellcolor{gray!50} &\cellcolor{gray!50} &\cellcolor{gray!50} &\cellcolor{gray!50} &\cellcolor{gray!50} &\cellcolor{gray!50} &\cellcolor{gray!50} &\cellcolor{gray!50} &\cellcolor{gray!50} &\cellcolor{gray!50} & \\
			\hline
			Finalização da documentação  & & & & & & & & & & & & & & & & &\cellcolor{gray!50} & \\
			\hline
			Defesa do TCC & & & & & & & & & & & & & & & & & &\cellcolor{gray!50} \\
			\hline
		\end{tabular}
		\vspace{0.5cm}	\\Fonte: Autoria própria (2023)
	\end{center}
\end{quadro}


\bibliography{referencias}
\vspace{1cm}

\begin{center}
	\begin{tabular}{|m{10cm}|m{5cm}|}
		\hline
		\textbf{Alunos} & \textbf{Assinatura} \\
		\hline
		 \textit{\imprimirprontuarioA - \imprimirautorA} &  \\
		\hline
		\textit{\imprimirprontuarioB - \imprimirautorB} &  \\
		\hline
		\textit{\imprimirprontuarioC - \imprimirautorC} & \\
		\hline
	\end{tabular}
\end{center}

\vspace{0.5cm}

\begin{center}
	\begin{tabular}{|m{10cm}|m{5cm}|}
		\hline
		\textbf{Orientador(es)} & \textbf{Assinatura} \\
		\hline
		\textit{\imprimirorient} & \\
		\hline
		\textit{\imprimircoorient} & \\
		\hline
	\end{tabular}
\end{center}

\end{document}

